\scanpage{136}\mspage{- 114 -}\hfill IX

qui est injective (puisque la suite exacte $0\longrightarrow V^*\longrightarrow K^*\longrightarrow\mathbb Z\longrightarrow0$ est scindée), donc elle est injective, ce qui implique (grâce au \unsure{sorite} 9.3) que si $X_\eta$ et $Y_\eta$ sont isomorphes, il en est de même de $X$ et de $Y$. On gagne donc par a).

Evidemment a) et b) impliquent que le foncteur $X\longmapsto X_\eta$ est \uline{pleinement} fidèle.

c)\quad Soit $X_\eta$ une extension panachée sur $\eta$, prouvons qu'elle se prolonge à $S$ si et seulement si les accouplements (9.4.6.2) sont nuls. Ici encore, l'unicité déjà prouvée montre que la question est locale pour la topologie étale, ce qui nous\unsure{,} ramène encore au cas $S$ strictement local. Alors la conclusion est vraie par construction, comme on avait signalé juste avant 9.4.6.

Ainsi la démonstration de 9.4.7 est achevée.

9.5.\quad Supposons maintenant fixé un nombre premier $\ell$, des pro-$\ell$-groupes de Barsotti-Tate sur $S$,
\[
 P=(P(n))_{n\geq0}\ ,\ R=(R(n))_{n\geq0}\ ,\quad Q=(Q(n))_{n\geq0}\ ,\tag*{(9.5.1)}
\]
et des extensions de pro-$\ell$-groupes de Barsotti-Tate
\[
 0\longrightarrow P\longrightarrow F\longrightarrow R\longrightarrow0\ ,\quad
 0\longrightarrow R\longrightarrow E\longrightarrow Q\longrightarrow0\ ,\tag*{(9.5.2)}
\]
définissant donc pour tout entier $n\geq0$ des extensions de $\Lambda_n$-Modules où $\Lambda_n=\mathbb Z/\ell^{n+1}\mathbb Z$ :
\[
 0\longrightarrow P(n)\longrightarrow F(n)\longrightarrow R(n)\longrightarrow0,\quad
 0\longrightarrow R(n)\longrightarrow E(n)\longrightarrow Q(n)\longrightarrow\unsure{0|C}\ .\tag*{(9.5.3)}
\]
Nous supposerons $Q$ de type étale et $P$ de type multiplicatif. Supposons
% archive sheet number: 135
% stamp: numbered label 432; the same number appears in pencil above the strip.
